\begin{figure*}[t]
\centering
\begin{tikzpicture}[x=1cm,y=1cm,font=\footnotesize,
  box/.style={draw=black!55,rounded corners=2pt,align=center,minimum height=1.0cm,text width=2.6cm,fill=white},
  flow/.style={-{Stealth[length=2mm]},thick,draw=blue!60!black},
  bgflow/.style={-{Stealth[length=2mm]},draw=black!65,dashed},
  lab/.style={font=\scriptsize,align=center,fill=white,inner sep=2pt}]
\node[box,text width=1.9cm] (sender) at (0,0) {付款方钱包\\\textbf{授权付款}};
\node[box,text width=3.3cm,fill=blue!4] (org) at (4,0) {消费裁决组织\\4名独立成员\\\textbf{3份匹配签名}};
\node[box,fill=blue!4] (wallet) at (8.8,0) {收款方钱包\\验证并记录TXCer\\\textbf{可发起后继付款}};
\node[box] (next) at (13,0) {后继付款\\相同或其他\\已注册组织};
\draw[flow] (sender) -- node[lab,above]{交易} (org);
\draw[flow] (org) -- node[lab,above]{输出 + TXCer} (wallet);
\draw[flow] (wallet) -- node[lab,above]{直接证据} (next);
\node[box,text width=3.3cm,fill=green!4] (comm) at (4,-3) {公共委员会\\\textbf{执行子付款}\\单次消费；记录直接缺口};
\node[box,fill=green!4] (resolve) at (8.8,-3) {等待资金来源\\来源到达后\\关闭直接义务};
\node[box,fill=orange!6] (repair) at (13,-3) {授权修复\\备付扣款 + 修订\\来源迟到：偿还原组织\\后台历史物化};
\draw[bgflow] (org) -- node[lab,right]{公共提交} (comm);
\node[lab,left] at (2.1,-1.25) {后台成员证据安装（INSTALL）\\在组织内部\\复制完整证据};
\draw[bgflow] (comm) -- node[lab,above]{直接\\义务} (resolve);
\draw[bgflow] (resolve) -- node[lab,above]{到期时\\仍未成立} (repair);
\node[align=center,font=\scriptsize,text=black!75] at (8.6,-4.5) {资金来源的落实保持已授权输出与后继引用不变。\\钱包与成员跟随经认证的公共区块更新状态；无需委员会逐笔通知。};
\end{tikzpicture}
\caption{Next中的连续支付与资金来源的公共落实。实线箭头表示收款人的续花路径，虚线箭头表示后台工作。每笔后继付款均需独立授权；尚未落实的资金来源由其直接发行组织承担责任，而不是作为等待祖先交易的要求向后传播。历史修复由后续公共命令授权；赔付后来源成功执行时偿还原赔付组织，不再向收款方重复入账。}
\label{fig:architecture}
\end{figure*}
